\scanpage{126}
\mspage{2}
% archive foliation: 125
 de schémas en groupes finis et plats sur $S$
\diagram{\begin{tikzcd}[column sep=small]
e \arrow[r] & H_{\mathrm{inf}} \arrow[r] & H \arrow[r] & H_{\text{sép}} \arrow[r] & e
\end{tikzcd}}
où $H_{\mathrm{inf}}$ est purement infinitésimal (i.e. la projection $H_{\mathrm{inf}}\to S$ est géométriquement injectif) et où $H_{\text{sép}}$ est non ramifié sur $S$. (Dans le cas d'un corps de base $k$, une telle suite exacte est déduite d'une suite exacte analogue pour un groupe algébrique fini \inserted{$G$} sur $k$; d'ailleurs si $k$ est parfait, cette suite splitte canoniquement, car $G_{\text{réd}}$ est \inserted{alors} un sous-groupe de $G$, isomorphe à $G_{\text{sép}}$ par la projection $G\quad G_{\text{sép}}$. Faire attention cependant que $G_{\text{réd}}$ va opérer sur $G_{\mathrm{inf}}$, on aura seulement un produit semi-direct). Enfin

\deleted{La condition (v) \unsure{est essentiellement provisoire}.}

(vi) $S$ est réduit.

Les conditions (v) et (vi) ont une sale gueule, et sont essentiellement provisoires. Elles servent à \deleted{\unsure{justifier la vérité}} \inserted{\unsure{cause}} du

\uline{Lemme}. Soit $H$ comme dans (v), $S$ comme dans (vi), $P$ un fibré principal homogène sous $H$, $Q$ un autre, $u,v$ deux isomorphismes de $P$ dans $Q$ transformant un \inserted{``}point \deleted{\unsure{géométrique}} \inserted{marqué''} donné en un autre donné, alors $u=v$. (On est ramené, en tordant $H$, au cas où $P$ est trivial, et alors cela résulte du fait suivant: une section de $Q$ est un isomorphisme de $S$ sur une composante connexe de $Q_{\text{réd}}$).

Remarque néanmoins que les conditions (i) à (vi) n'excluent pas des groupes structuraux tordus (relativement à un corps de base $k$).

Soit maintenant $a$ un ``point marqué'' de $S$ (i.e. une extension algébriquement close du corps résiduel d'un $s\in S$). Pour tout $G\in\underline G$, on désigne par $Z(S,a;G)$ ou simplement $Z(G)$ l'ensemble des classes (à isomorphisme près) de fibrés principaux homogènes sur $S$, de groupe $F(G)$, munis d'un point marqué au dessus de $a$. Evidemment $Z(G)$ est un foncteur de $\underline G$ dans la catégorie des ensembles. Ce foncteur satisfait aux conditions suivantes: 1) Il commute aux produits (\deleted{trivial) \unsure{et}} car $F$ y commute) 2) il commute aux noyaux de couples, en d'autres termes si
\diagram{\begin{tikzcd}[column sep=small]
G'' \arrow[r] & G \arrow[r,shift left=.6ex] \arrow[r,shift right=.6ex] & G'
\end{tikzcd}}
est exact dans $\underline G$, alors
\diagram{\begin{tikzcd}[column sep=small]
Z(G'') \arrow[r] & Z(G) \arrow[r,shift left=.6ex] \arrow[r,shift right=.6ex] & Z(G')
\end{tikzcd}}
est exacte, i.e. $Z(G'')$ s'identifie à l'ensemble des éléments de $Z(G)$ dont l'image dans $Z(G')$ par $Z(u)$ et $Z(v)$ est la même. (Ceci résulte de l'exactitude de $F$, et du lemme).
\figure{\begin{tikzpicture}[scale=.4]
\draw[line width=.6pt] (0,.1)--(.08,.32)--(.22,.35)--(.4,.53)--(.13,.42)--(.04,.3)--(.1,-.1)--(.21,-.22)--(.41,-.14)--(.23,-.3)--(.12,-.1)--(-.03,.17)--(.22,.23);
\end{tikzpicture}}
